\section{Технический проект}
\subsection{Общая характеристика организации решения задачи}

Разрабатывается серверная часть веб-приложения для чат-платформы, обеспечивающая API для взаимодействия между клиентской частью и базой данных. Сервер реализует следующие основные функции:
\begin{itemize}
	\item аутентификация и авторизация пользователей;
	\item обмен сообщениями в групповых чатах;
	\item личная переписка между пользователями;
	\item управление группами и правами участников;
	\item работа с вложениями и медиафайлами.
\end{itemize}

\subsection{Обоснование выбора технологии проектирования}

Для реализации серверной части выбраны следующие технологии:

\subsubsection{Язык программирования Python}

Python выбран благодаря:
\begin{itemize}
	\item простоте и читаемости кода;
	\item богатой экосистеме веб-фреймворков;
	\item хорошей поддержке работы с базами данных;
	\item кроссплатформенности.
\end{itemize}

\subsubsection{Интерфейс WSGI}

В качестве основы сервера используется WSGI (Web Server Gateway Interface) - стандартный интерфейс между веб-сервером и Python-приложениями. Преимущества:
\begin{itemize}
	\item легковесность;
	\item простота развертывания;
	\item совместимость с различными серверами (Waitress, Gunicorn и др.).
\end{itemize}

\subsubsection{База данных SQLite}

SQLite выбрана как:
\begin{itemize}
	\item встроенное решение, не требующее отдельного сервера;
	\item простое в настройке и использовании.
\end{itemize}

\subsubsection{Шифрование сообщений}

Рассмотрим пошагово, как происходит подготовка и шифрование текстового сообщения, отправляемого из клиентского приложения на сервер:

\begin{enumerate}
	\item формирование исходного текста:  
	пользователь вводит строку "Hello World". Эта строка подлежит шифрованию с использованием блочного алгоритма симметричного шифрования AES в режиме ECB. Поскольку алгоритм работает с блоками длиной 16 байт, исходное сообщение нужно дополнить до нужной длины;
	
	\item добавление выравнивающих байтов (padding):  
	длина строки "Hello World" составляет 11 байт. Согласно стандарту PKCS7, добавляются 5 служебных байтов со значением \textbackslash x05, чтобы итоговая длина составила 16 байт. После выравнивания строка имеет вид:  
	"Hello World\textbackslash x05\textbackslash x05\textbackslash x05\textbackslash x05\textbackslash x05";
	
	\item разбивка на блоки:  
	полученное сообщение разбивается на части по 16 байт. Если после паддинга длина становится кратна 16, сообщение состоит из одного блока; при большей длине — из нескольких. В данном примере:
	\begin{itemize}
		\item первый блок: Hello World\textbackslash x05\textbackslash x05
		\item второй блок (если он образован при добавлении нулей после паддинга): \textbackslash x05\textbackslash x05\textbackslash x05\textbackslash x05\textbackslash x05 и дополнительные байты
	\end{itemize}
	
	\item шифрование каждого блока:  
	каждый блок последовательно шифруется сессионным ключом. В результате формируются зашифрованные блоки байтов, например:
	\begin{itemize}
		\item первый блок превращается в A3F1...;
		\item второй блок — в 9B02...
	\end{itemize}
	
	\item объединение зашифрованных блоков:  
	зашифрованные блоки объединяются в одно сообщение. Финальный результат:  
	A3F1... + 9B02... = A3F19B02...
\end{enumerate}

На рисунке \ref{shifr} наглядно представлен пример шифрования текстового сообщения, отправляемого из клиентского приложения на сервер.

	\begin{figure}[H]
	\centering
	\includegraphics[width=0.9\linewidth]{"images/Пример шифрования"}
	\caption{Пример шифрования текстового сообщения}
	\label{shifr}
\end{figure}

Таким образом, даже короткая текстовая строка проходит сквозь цепочку выравнивания, блочной обработки и шифрования, прежде чем попасть на сервер в защищённом виде.

\subsubsection{Пример расшифровки сообщения}

Рассмотрим пошагово, как происходит расшифровка зашифрованного сообщения, полученного сервером:

\begin{itemize}
	\item получение зашифрованного сообщения. сервер получает конкатенированный поток зашифрованных блоков, например, «A3F19B02...»;
	\item разбивка зашифрованного сообщения на блоки фиксированной длины. поскольку алгоритм блочного шифрования работает с блоками по 16 байт, зашифрованное сообщение делится на отдельные блоки;
	\item расшифровка каждого блока. каждый блок обрабатывается алгоритмом AES (в режиме ECB) с использованием того же сессионного ключа, который применялся при шифровании, в результате чего зашифрованные блоки преобразуются обратно в исходные блоки, содержащие текст с выравниванием;
	\item конкатенация расшифрованных блоков. полученные блоки объединяются в одно сообщение, содержащее исходный текст вместе с добавленными служебными байтами (padding);
	\item удаление padding по стандарту PKCS\#7. из объединенного сообщения удаляются дополнительные байты, что позволяет восстановить исходное сообщение, например, «Hello World».
\end{itemize}

На рисунке \ref{deshifr} наглядно представлен пример расшифровки текстового сообщения.

\begin{figure}[H]
	\centering
	\includegraphics[width=0.9\linewidth]{"images/Пример расшифровки"}
	\caption{Пример расшифровки текстового сообщения}
	\label{deshifr}
\end{figure}

Таким образом, процесс расшифровки обеспечивает восстановление оригинального текста сообщения, гарантируя передачу данных в защищённом виде и корректное их восстановление на стороне сервера.

\subsection{Диаграмма компонентов серверной части}

На рисунке \ref{fig:-components} представлена архитектура серверной части приложения.

\newpage
\begin{landscape}
	\begin{figure}[ht]
		\centering
		\includegraphics[width=0.9\linewidth]{"images/Диаграмма компонентов"}
		\caption{Диаграмма компонентов серверной части}
		\label{fig:-components}
	\end{figure}
\end{landscape}

\subsubsection{WSGI-приложение (app.py)}
Данный модуль является точкой входа в серверное приложение. Основная функция app(environ, start\_response) реализует следующие задачи:
\begin{itemize}
	\item принимает входящие HTTP-запросы (содержащие URL, параметры запроса и другие данные) через окружение environ и функцию start\_response;
	\item определяет маршрутизацию запроса, обращаясь к словарю маршрутов из модуля routes.py;
	\item обрабатывает возникшие исключения и в случае ошибок перенаправляет обработку на соответствующие view (например, InternalServerErrorView);
	\item интегрируется с сервером Waitress для развертывания WSGI-приложения в производственной среде.
\end{itemize}

\subsubsection{Маршрутизация (routes.py)}
Модуль маршрутизации содержит словарь, в котором URL-паттерны сопоставляются с соответствующими классами-обработчиками (view). Его ключевые функции:
\begin{itemize}
	\item функция route(url) проходит по ключам словаря и на основе регулярных выражений определяет, какой view-класс следует вызвать для обработки запроса;
	\item обеспечивает централизованное распределение входящих запросов между различными модулями, например, для авторизации, работы с сообщениями, управления группами и обработки ошибок.
\end{itemize}

\subsubsection{Модуль MIME-типов (mimes.py)}
Этот модуль отвечает за определение MIME-типов файлов, отдаваемых сервером. Основные аспекты:
\begin{itemize}
	\item содержит словарь, сопоставляющий расширения файлов с соответствующими MIME-типами;
	\item функция get\_mime(file\_name) ищет расширение в имени файла и возвращает корректный MIME-тип, что критически важно для правильного отображения статических ресурсов (HTML, CSS, JS, изображения и т.д.).
\end{itemize}

\subsubsection{Инициализация базы данных}
На этапе старта приложения вызывается функция initialize\_database(), которая обеспечивает следующее:
\begin{itemize}
	\item создание необходимых таблиц (таких как users, groups, attachments, group\_members, group\_messages, private\_messages, session) в базе данных SQLite;
	\item создание общего чата (например, с названием «Общий чат») и установка индексирования для ускорения поиска по сообщениям;
	\item обеспечение корректной инициализации структур данных, необходимых для работы всего приложения.
\end{itemize}

\subsubsection{Обработка статических файлов (serve\_static)}
Функция serve\_static(environ, start\_response) реализует следующую логику:
\begin{itemize}
	\item по заданному URL извлекается путь к запрашиваемому статическому файлу;
	\item если файл найден, определяется его MIME-тип с использованием get\_mime и отдается содержимое файла;
	\item в случае отсутствия файла возвращается ответ с ошибкой 404 и соответствующим типом контента (text/plain).
\end{itemize}

\subsubsection{View-компоненты (пакет views)}
Пакет views содержит набор классов для обработки HTTP-запросов, каждая категория которых отвечает за определённый функционал:
\begin{itemize}
	\item базовые классы View и TemplateView (модуль base.py) реализуют общую логику формирования ответов, работы с шаблонами и обработки ошибок;
	\item модули в auth.py (например, LoginView, RegisterView, LogoutView) обрабатывают регистрацию, логин, создание сессий и выход пользователей;
	\item модули в message.py реализуют функции получения, отправки, редактирования и удаления сообщений для общего и группового чатов;
	\item p\_chat.py отвечает за получение списка личных чатов и проверку их обновлений;
	\item модули в search.py позволяют выполнять поиск пользователей и сообщений с различными параметрами;
	\item модуль users.py реализует функции получения данных о пользователе, списке членов общего чата и изменение ролей.
\end{itemize}

\subsubsection{Модели данных (пакет models)}
Пакет models отвечает за взаимодействие с базой данных и бизнес-логику работы с данными:
\begin{itemize}
	\item модуль usermodel.py реализует операции создания, аутентификации и получения данных пользователей;
	\item модуль messagemodels.py (и связанные с ним модули) осуществляют CRUD-операции для сообщений в общем, групповых и приватных чатах, а также поддерживают обработку вложений;
	\item модуль groupmodels.py отвечает за создание групп, управление участниками, изменение ролей и переименование групп;
	\item модуль session.py реализует создание, получение и удаление сессионных данных, а также генерацию криптографически стойких ключей.
\end{itemize}

\subsubsection{Утилиты и вспомогательные функции (utils)}
В данном модуле находятся вспомогательные функции, обеспечивающие единообразие и упрощение логики обработки:
\begin{itemize}
	\item функция json\_response стандартизирует формат ответов в формате JSON;
	\item функция forbidden\_response формирует ответы для случаев отсутствия прав доступа;
	\item реализованы дополнительные функции для работы с MIME-типами, чтения файлов, логирования и обработки исключений.
\end{itemize}

\subsubsection{HTTP-сервер (Waitress)}
Для развертывания серверного приложения используется Waitress:
\begin{itemize}
	\item waitress запускает WSGI-приложение, обеспечивая его работу в условиях продакшн-среды;
	\item конфигурация сервера производится в файле run.py посредством вызова waitress.serve(app, host='0.0.0.0', port=8000);
	\item сервер принимает HTTP-запросы от клиентов, передавая их в приложение для дальнейшей обработки.
\end{itemize}

Структура серверной части корпоративного мессенджера организована модульно и охватывает все ключевые аспекты работы приложения — от маршрутизации и обработки HTTP-запросов до управления данными пользователей, сообщений, сессий и вложений. Такое разделение компонентов обеспечивает читаемость, масштабируемость и удобство сопровождения кода. Централизованная логика взаимодействия с базой данных, модульная архитектура view-обработчиков, отдельные модели и сервисные утилиты позволяют гибко развивать функциональность, не нарушая стабильность уже реализованных компонентов. Серверная часть эффективно взаимодействует с клиентской и модельной, играя роль устойчивого посредника между интерфейсом пользователя и слоем хранения данных.

\subsection{Основные классы и функции}

Модульная организация кода классов представлений и вспомогательных функций. Классы и функции распределены по специализированным модулям внутри директории views (а также в отдельных модулях для утилит). Такой подход позволяет обеспечить более удобное чтение кода, модульное тестирование и масштабируемость системы.

\subsubsection{Базовый модуль base.py} 
Базовый модуль, реализующий общий функционал для всех представлений, включает:
\begin{itemize}
	\item \textbf{View} класс для обработки базовых HTTP-запросов и формирования стандартных ответов;
	\item \textbf{TemplateView} класс для работы с файловыми шаблонами и генерации HTML-страниц;
	\item \textbf{IndexView} класс для отображения главной страницы приложения.
\end{itemize}

Данный модуль отвечает за:

\begin{itemize}
	\item обработку HTTP-запросов;
	\item чтение и интерпретацию файлов шаблонов;
	\item формирование ответов сервера.
\end{itemize}

\subsubsection{Модуль auth.py}
В данном модуле размещены классы, связанные с работой пользователей:
\begin{itemize}
	\item \textbf{RegisterView} класс для обработки регистрации новых пользователей;
	\item \textbf{LoginView} класс для обработки аутентификации и входа в систему.
\end{itemize}

\subsubsection{Модуль group.py}
В данном модуле размещены классы, связанные с работой групп:
\begin{itemize}
	\item \textbf{GetGroupMembersView} класс для получения списка участников группы;
	\item \textbf{LeaveGroupView} класс для выхода пользователя из группы;
	\item \textbf{CreateGroupView} класс для создания новой группы;
	\item \textbf{AddToGroupView} класс для добавления пользователей в группу;
	\item \textbf{GetGroupView} класс для получения подробной информации о группе;
	\item \textbf{CheckGroupAccessView} класс для проверки прав доступа пользователя в группе;
	\item \textbf{GetGroupNameView} класс для получения названия группы по её идентификатору;
	\item \textbf{CheckGroupsUpdateView} класс для проверки обновлений в информации о группах;
	\item \textbf{RenameGroupView} класс для переименования группы;
	\item \textbf{RemoveFromGroupView} класс для удаления участников из группы.
\end{itemize}

Данный модуль отвечает за:

\begin{itemize}
	\item управление данными групп пользователей;
	\item контроль доступа и администрирование групп;
	\item обеспечение актуальности информации о группах.
\end{itemize}

\subsubsection{Модуль error.py}
В данном модуле размещены классы, связанные с обработкой ошибок:
\begin{itemize}
	\item \textbf{NotFoundView} класс для отображения ошибки 404 (Не найдено);
	\item \textbf{InternalServerErrorView} класс для обработки ошибки 500 (Внутренняя ошибка сервера);
	\item \textbf{ForbiddenView} класс для генерации ответа 403 (Доступ запрещён).
\end{itemize}

Данный модуль отвечает за:

\begin{itemize}
	\item обработку исключительных ситуаций и ошибок;
	\item вывод информативных сообщений при сбоях;
	\item защиту приложения от некорректных запросов.
\end{itemize}

\subsubsection{Модуль message.py}
В данном модуле размещены классы, связанные с обработкой сообщений:
\begin{itemize}
	\item \textbf{GetMessageView} класс для получения списка сообщений;
	\item \textbf{SendMessageView} класс для отправки новых сообщений;
	\item \textbf{DeleteMessageView} класс для удаления сообщений по идентификатору;
	\item \textbf{EditMessageView} класс для редактирования ранее отправленных сообщений;
	\item \textbf{GetGroupMessageView} класс для получения сообщений группового чата;
	\item \textbf{GetPrivateMessageView} класс для получения личных сообщений;
	\item \textbf{SendPrivateMessageView} класс для отправки личных сообщений;
	\item \textbf{SendSystemMessageView} класс для отправки системных уведомлений;
	\item \textbf{CheckMessagesView} класс для проверки наличия новых сообщений;
	\item \textbf{CheckEditedMessageView} класс для проверки изменений в сообщениях.
\end{itemize}

Данный модуль отвечает за:

\begin{itemize}
	\item организацию процесса обмена сообщениями между несколькими пользователями;
	\item поддержание целостности переписок;
	\item обеспечение возможности модификации и удаления сообщений.
\end{itemize}

\subsubsection{Модуль pchat.py}
В данном модуле размещены классы, связанные с обработкой личных чатов:
\begin{itemize}
	\item \textbf{GetPrivateChatsView} класс для получения списка личных чатов пользователя;
	\item \textbf{CheckPrivateChatsUpdateView} класс для проверки обновлений в личных чатах.
\end{itemize}

Данный модуль отвечает за:

\begin{itemize}
	\item организацию и управление личными чатами;
	\item обеспечение оперативного доступа к истории переписки.
\end{itemize}

\subsubsection{Модуль search.py}
В данном модуле размещены классы, связанные с поиском:
\begin{itemize}
	\item \textbf{SearchUsersView} класс для поиска пользователей по заданным критериям (например, по имени или email);
	\item \textbf{SearchMessagesView} класс для поиска сообщений по содержимому и дате.
\end{itemize}

Данный модуль отвечает за:
	
\begin{itemize}
	\item обеспечение функционала поиска по базе данных пользователей и сообщений;
	\item быстрое извлечение релевантных данных;
	\item улучшение навигации пользователя по приложению.
\end{itemize}

\subsubsection{Модуль users.py}
В данном модуле размещены классы, связанные с обработкой данных пользователей:
\begin{itemize}
	\item \textbf{GetUserIDView} класс для получения идентификатора текущего пользователя;
	\item \textbf{GetGeneralMembersView} класс для получения списка всех пользователей системы;
	\item \textbf{ChangeMemberRoleView} класс для изменения роли участника в группе.
\end{itemize}

Данный модуль отвечает за:

\begin{itemize}
	\item управление информацией о пользователях;
	\item администрирование и обновление пользовательских данных;
	\item поддержание актуальности и целостности данных системы.
\end{itemize}


\subsubsection{Вспомогательные модули.}
Для поддержки работы основных классов выделены несколько вспомогательных модулей: \texttt{routes}, \texttt{mimes}, \texttt{utils}, \texttt{app} и \texttt{run}, каждый из которых реализует специфичный набор функций:
\begin{itemize}
	\item \textbf{session} обеспечивает генерацию ключей для шифрования сообщений с помощью AES;
	\item \textbf{routes} обеспечивает маршрутизацию HTTP-запросов к соответствующим представлениям, что позволяет определить, какой модуль и функция должны обработать тот или иной запрос;
	\item \textbf{mimes} реализует функции для определения MIME-типов файлов на основе расширений или содержимого, обеспечивая корректное формирование заголовков ответов;
	\item \textbf{utils} содержит набор вспомогательных утилит, включая функции для генерации JSON-ответов, обработки ошибок, а также хеширования и проверки паролей;
	\item \textbf{app} отвечает за инициализацию и конфигурацию серверного приложения, настройку параметров окружения и подключение к базовым сервисам, так же и к базе данных);
	\item \textbf{run} предназначен для управления запуском и жизненным циклом серверного приложения, включая старт, остановку и мониторинг его работы.
\end{itemize}


\subsubsection{Вспомогательные функции.}
Для поддержки работы основных классов выделен отдельный модуль \texttt{utils}, содержащий функции:
\begin{itemize}
	\item \textbf{json\_response} формирование JSON-ответов для клиентской части;
	\item \textbf{forbidden\_response} генерация ответа 403 Forbidden при отсутствии прав доступа;
	\item \textbf{get\_mime} определение MIME-типа для корректной передачи файлов;
	\item \textbf{check\_group\_permissions} проверка прав доступа пользователя в группе по заданным параметрам;
	\item \textbf{hash\_password} и \textbf{check\_password} функции для безопасного хеширования паролей и их верификации.
\end{itemize}

